\section{Introduction}\label{sec:intro}
\sectionstate{partial}{first prose draft 2026-11-18; \texttt{RECONSTRUCTION\_BASIS.md}; \texttt{WAVEFORM\_FIRST\_THREADING.md}}
A Micromegas with a drift gap of a few centimetres can be read as a small time projection chamber: the electrons
of an inclined track reach the mesh one after the other, the strip they land on gives one coordinate and their
arrival time the depth, and a single plane measures a track segment with its angle (the micro-TPC
mode~\cite{paperI}). The usual reconstruction reduces each strip to one hit, a time and an amplitude, and fits a
line through the hit times. That line is biased. A strip collects charge from several drift depths, because the
charge spreads laterally on its way to the readout, and the one time it reports is an aggregate of them; the
resulting ladder of hit times is compressed and the track reads too steep, by an amount that depends on the angle and
that no choice of threshold or time estimator removes. The measurement of that bias, its scaling law and what a
hit-based readout would have to record to avoid it are the subject of a companion paper~\cite{paperIIb}.

This paper gives the reconstruction that avoids the bias when the waveforms are available. The lateral spreads that
mix the depths are few and physical: drift diffusion, the prompt footprint of the avalanche and its induction, and,
on a resistive anode, the spreading of the charge along the resistive strips. We put them in a forward model of the
sampled strip signals of a straight track and fit that model to the whole picture of amplitude against strip and
time, instead of fitting a line to times extracted strip by strip. What is given up is generality: the model
describes one straight track per fit, not the curved or scattered trajectories a true TPC can follow.

The detector is the MX17 chamber of~\cite{paperI}: a resistive-strip bulk Micromegas with a 30\,mm drift gap, two
strip views at 0.78\,mm pitch, read by the DREAM front end, which samples every strip every 60\,ns. Five chambers
(A to E) were measured with cosmic muons on a test bench against an external tracker (M3), and chamber E in the SPS H4
beam, where it was also rotated to a known angle. Section~\ref{sec:kernel} measures how charge reaches the
neighbouring strips; \cref{sec:hitsfail} shows what this does to hit times; \cref{sec:model} defines the model, the fit
and the calibration, which uses the chamber's own data only; \cref{sec:floor} compares the resolution with what the
ionisation and diffusion allow; \cref{sec:tiers} compares, on the same events, the hit line, the same model fitted
to hit lists, and the waveform fit; \cref{sec:fleet} gives the performance of the five chambers and the test at a
known angle.
